\section{La solution aqueuse d'acide méthanoïque}
L'acide méthanoïque $\ce{HCOOH}$, couramment appelé acide formique, est un
liquide piquant et corrosif qui existe à l'état naturel dans l'organisme des
fourmis rouges.

Cet exercice vise~:
\begin{itemize}
    \item l'étude d'une solution aqueuse d'acide méthanoïque~;
    \item le dosage d'une solution aqueuse d'acide méthanoïque~;
    \item la comparaison du comportement de deux acides.
\end{itemize}

\subsection{Étude de la solution aqueuse d'acide méthanoïque}
On dispose d'une solution aqueuse $\mathrm{(S_{A})}$ d'acide méthanoïque
$\ce{HCOOH_{(aq)}}$ de volume $V=\qty{1}{\liter}$, de concentration molaire
$C_{A}=\qty{0,10}{\M}$ et de $\mathrm{pH}=2,4$.
\begin{enumerate}
    \item Définir un acide selon Brönsted.
    \item Écrire l'équation modélisant la transformation chimique entre l'acide
          méthanoïque et l'eau.
    \item Recopier sur votre copie le tableau d'avancement et le compléter.
          \begin{table}[H]
          \renewcommand{\arraystretch}{1.5}
          \begin{tabularx}{\textwidth}{|
              >{\centering\arraybackslash}m{2.5cm}|
              >{\centering\arraybackslash}m{3.5cm}|C|C|C|C|}
              \hline
              \multicolumn{2}{|c|}{Équation chimique} &
              \multicolumn{4}{c|}{\dotfill} \\
              \hline
              État du système & Avancement de la réaction en (\unit{\mol}) &
              \multicolumn{4}{c|}{Quantité de matière en (\unit{\mol})} \\
              \hline
              État initial & 0 & \dotfill & \dotfill & \dotfill & \dotfill \\
              \hline
              État intermédiaire & $x$ & \dotfill & \dotfill & \dotfill &
              \dotfill \\
              \hline
              État final & $x_{f}$ & \dotfill & \dotfill & \dotfill &
              \dotfill \\
              \hline
          \end{tabularx}
          \end{table}
    \item Calculer la valeur de l'avancement final $x_{f}$ de cette réaction.
    \item Calculer le taux d'avancement final $\tau$ de cette réaction.
          Conclure.
    \item Montrer que le quotient de réaction à l'état d'équilibre du système
          chimique s'écrit~:
          $Q_{r,\eq}=\frac{10^{-2\cdot\mathrm{pH}}}{C_{A}-10^{-\mathrm{pH}}}$.
          Calculer sa valeur.
    \item Déduire la valeur de la constante d'équilibre $\mathrm{K}$ associé à
          l'équation de la réaction. 
\end{enumerate}




\subsection{Dosage de la solution aqueuse d'acide méthanoïque}
\begin{minipage}{0.74\linewidth}
À fin de vérifier la valeur de la concentration molaire $C_{A}$ de la solution
$\mathrm{(S_{A})}$, on réalise un titrage acido-basique.

Dans un bécher, on verse un volume $V_{A}=\qty{20,0}{\mL}$ de cette solution et
on y ajoute progressivement une solution aqueuse $\mathrm{(S_{B})}$ d'hydroxyde
de sodium $\ce{Na^+_{(aq)} + HO^-_{(aq)}}$ de concentration molaire
$C_{B}=\qty{0,25}{\M}$. Les coordonnées du point d'équivalence sont~:
$(V_{B,\mathrm{E}}=\qty{8,0}{\mL}~; \mathrm{pH_{E}}=8,2)$.

Le montage expérimental utilisé pour réaliser ce dosage est représenté sur
la figure ci-contre.
\end{minipage}
\hfill
\begin{minipage}{0.24\linewidth}
    \flushright
    \begin{tikzpicture}
        \node {\includegraphics[width=\textwidth]{2025_03_31_52b4902f5779608975f8g-2}};
    \end{tikzpicture}
\end{minipage}
\begin{enumerate}
    \item Nommer les éléments correspondants aux numéros indiqués sur le montage
          de la figure.
    \item Écrire l'équation de la réaction qui se produit entre l'acide
          méthanoïque $\ce{HCOOH_{(aq)}}$ et les ions hydroxydes
          $\ce{HO^-_{(aq)}}$ au cours du dosage, sachant qu'elle est totale.
    \item Vérifier la valeur de $C_{A}$.
    \item Parmi les deux indicateurs colorés suivants, quel est celui qui
          convient le mieux à ce dosage~? Justifier.

          \begin{minipage}{\linewidth}
              \begin{table}[H]
                  \centering
                  \renewcommand{\arraystretch}{1.5}
                  \begin{tabularx}{\textwidth}{|p{4cm}|C|C|C|}
                  \hline
                  \textbf{Indicateur coloré} & \textbf{Teinte acide} &
                  \textbf{Zone de virage} & \textbf{Teinte basique} \\
                  \hline
                  Rouge de crésol & Jaune & $7,2-8,8$ & Rouge \\
                  \hline
                  Alizarine & Rouge & $11,0-12,4$ & Violet \\
                  \hline
                  \end{tabularx}
              \end{table}
          \end{minipage}
    \item Pour un volume versé $V_{B}=\frac{V_{B,E}}{2}$ de la solution
          $\mathrm{(S_{B})}$, le $\mathrm{pH}$ du mélange dans le bécher vaut
          $\mathrm{pH}=3,8$ et $[\ce{HCOOH_{(aq)}}]=[\ce{HCOO^-_{(aq)}}]$.

          Calculer la constante d'acidité $\mathrm{K_{A}}$ du couple
          $(\ce{HCOOH_{(aq)}}\ttslash{}\ce{HCOO^-_{(aq)}})$.
\end{enumerate}




\subsection{Comportement de deux acides en solution aqueuse}
On considère une seconde solution aqueuse $\mathrm{(S^{\prime})}$ d'acide
propanoïque $\ce{C2H5COOH}$ de concentration molaire
$C_{A}^{\prime}=\qty{0,10}{\M}$. La valeur du taux d'avancement final de la
réaction de l'acide propanoïque avec l'eau est $\tau^{\prime}=\num{1,16e-3}$.
\begin{enumerate}
    \item En comparant $\tau^{\prime}$ avec $\tau$ le taux d'avancement final de
          la réaction d'acide méthanoïque avec l'eau, indiquer lequel des deux
          acides est le plus dissocié en solution.
    \item Comparer les constantes d'acidité
          $\mathrm{K_{A}}(\ce{HCOOH_{(aq)}}/\ce{HCOO^-_{(aq)}})$ et
          $\mathrm{K_{A}}(\ce{C2H5COOH_{(aq)}}/\ce{C2H5COO^-_{(aq)}})$.
\end{enumerate}


